\section{Differentiate one scalar program before a tensor network}
Chapter 2 separated a declared objective from evidence that it measures the
intended capability. Now freeze the objective and ask a local question: which
operation transmits a change in one parameter to the loss?

Begin with two scalars, not an automatic differentiation interface:
\begin{equation}
x=2,\qquad w=3,\qquad a=wx,\qquad L=\tfrac12a^2.
\end{equation}
The forward values are $a=6$ and $L=18$. At the multiply, either input can
change the result. At the square, the local rate of loss change is $a$.
The reverse calculation must follow those dependencies, not the order in
which variable names happen to appear in a source file.

\NativeFigure{scalar-pullback}{Solid arrows carry forward values. Dashed arrows
carry reverse sensitivities. Both inputs enter multiplication independently;
the reverse pass produces a separate derivative for each input.}{fig:textbook-pullback}

Write a bar over a quantity to denote the derivative of the final loss with
respect to it. Starting from $\bar L=1$, the chain rule gives
\begin{align}
\bar a&=\bar L\,a=6,\\
\bar w&=\bar a\,x=12,\\
\bar x&=\bar a\,w=18.
\end{align}
These are not three gradients that should be averaged. They are derivatives
with respect to three different variables. The input $x$ can be fixed data
while its mathematical derivative is still well defined. Computing $\bar x$
does not imply that the optimizer updates the data.

\textbf{Prediction problem.} Update only $w$ with learning rate $.01$.
What are the new weight and loss for the same fixed input?

\textbf{Worked answer.} The update is $w'=3-.01(12)=2.88$.
The new loss is $\frac12(2\cdot2.88)^2=16.5888$. The loss decreases in this
particular smooth example. A computed gradient does not guarantee improvement
for an arbitrary finite learning rate or on a different evaluation dataset.

We can check the derivative without reusing the chain-rule program:
\begin{equation}
\frac{L(w+h)-L(w-h)}{2h}
=\frac{2(w+h)^2-2(w-h)^2}{2h}=4w=12.
\end{equation}
The equality is exact in real arithmetic for this quadratic and nonzero $h$.
Floating-point evaluation introduces subtraction and rounding error. Both
evaluations must use the same input and state; changing either means we are
subtracting values of different functions.

The matrix derivation that follows is this same calculation with shared
parameters, indexed sums and explicit shapes. A broadcast bias receives
contributions from several rows, so its reverse rule sums those contributions.
A nonlinear operation multiplies by its local derivative. Tensor notation
compresses these relationships; it does not replace their explanation.
